\sect{विनायक-व्रत-कल्पः}
\centerline{\small{(मूलम्—श्रीव्रतरत्नाकरः, प्रथमो भागः)}}

\twolineshloka
{आसीत्पुरा चन्द्रवंशे राजा धर्म इति श्रुतः}
{स्वराज्ये दैवयोगेन ज्ञातिभिः कुटिलैर्हृते} %॥१॥

\twolineshloka
{अनुजैर्भार्यया सार्थं जगाम गहनं वनम्}
{बहुवृक्षसमाकीर्णं नानामृगसमन्वितम्} %॥२॥

\twolineshloka
{बहुपक्षिकुलोपेतं व्याघ्रभल्लूकसङ्कुलम्}
{तत्र तत्र समाविष्टा मुनयो ब्रह्मवादिनः} %॥३॥

\twolineshloka
{आदित्यसन्निभाः सर्वे सर्वे वह्निसमप्रभाः}
{तेजोमण्डलसङ्काशा वायुपर्णाम्बुभक्षकाः} %॥४॥

\twolineshloka
{अग्निहोत्ररता नित्यम् अतिथीनां च पूजकाः}
{ऊर्ध्वबाहवो निरालम्बाः सर्वे मुनिगणास्तथा} %॥५॥

\twolineshloka
{तान् पश्यन् धर्मराजोऽपि सम्भ्रमेण समन्वितः}
{सूताश्रमं समासाद्य सूतं दृष्ट्वा ससम्भ्रमः} %॥६॥

\onelineshloka
{नत्वा च भार्यया सार्थम् अनुजैः समुपाविशत्} %॥७॥

\uvacha{धर्म उवाच}

\twolineshloka
{सूत सूत महाप्राज्ञ सर्वशास्त्रविशारद}
{वयं च भार्यया सार्थं ज्ञातिभिः परिपीडिताः} %॥८॥

\twolineshloka
{स्वराज्यं सकलं चैव पुत्राश्चापहृता हि नः}
{तव दर्शनमात्रेण सर्वं दुःखं विनाशितम्} %॥९॥

\onelineshloka
{ममोपरि कृपां कृत्वा व्रतं ब्रूहि दयानिधे} %॥१०॥

\uvacha{सूत उवाच}

\twolineshloka
{व्रतं सम्पत्करं नृणां सर्वसौख्यप्रवर्धनम्}
{शृणुध्वं पाण्डवाः सर्वे व्रतानामुत्तमं व्रतम्} %॥११॥

\twolineshloka
{रहस्यं सर्वपापघ्नं पुत्रपौत्राभिवर्धनम्}
{व्रतं साम्बशिवेनैव स्कन्दस्योद्बोधितं पुरा} %॥१२॥

\twolineshloka
{कैलासशिखरे रम्ये नानामुनिनिषेविते}
{मन्दारविटपिप्रान्ते नानामणिविभूषिते} %॥१३॥

\twolineshloka
{हेमसिंहासनासीनं शङ्करं लोकशङ्करम्}
{पप्रच्छ षण्मुखस्तुष्टो लोकानुग्रहकाङ्क्षया} %॥१४॥

\uvacha{स्कन्द उवाच}

\twolineshloka
{केन व्रतेन भगवन् सौभाग्यम् अतुलं भवेत्}
{पुत्रपौत्रान् धनं लब्ध्वा मनुजः सुखमेधते} %॥१५॥

\onelineshloka
{तन्मे वद महादेव व्रतानामुत्तमं व्रतम्} %॥१६॥

\uvacha{ईश्वर उवाच}

\twolineshloka
{अस्ति चात्र महाभाग गणनाथप्रपूजनम्}
{सर्वसम्पत्करं श्रेष्ठम् आयुष्कामार्थसिद्धिदम्} %॥१७॥

\twolineshloka
{मासि भाद्रपदे शुक्ल-चतुर्थ्यां व्रतमाचरेत्}
{प्रातः स्नात्वा शुचिर्भूत्वा नित्यकर्म समाचरेत्} %॥१८॥

\twolineshloka
{स्वशक्त्या गणनाथस्य स्वर्णरौप्यमयाकृतिम्}
{अथवा मृण्मयं कुर्याद् वित्तशाठ्यं न कारयेत्} %॥१९॥

\twolineshloka
{स्वगृहस्योत्तरे देशे मण्डपं कारयेत्ततः}
{तन्मध्ये अष्टदलं पद्मं यवैर्वा तण्डुलेन वा} %॥२०॥

\twolineshloka
{प्रतिमां तत्र संस्थाप्य पूजयित्वा प्रयत्नतः}
{श्वेतगन्धाक्षतैः पुष्पैः बहुविधैश्च वारिभिः} %॥२१॥

\twolineshloka
{धूपैर्दीपैश्च नैवेद्यैः मोदकैर्घृतपाचितैः}
{एकविंशतिसंख्यानि नारिकेलफलान्यपि} %॥२२॥

\twolineshloka
{रम्भाजम्बुकपित्थौघान् इक्षुखण्डांश्च तावतः}
{एवमन्यफलापूपैर्नैवेद्यं कारयेत्सुत} %॥२३॥

\twolineshloka
{नृत्यगीतैश्च वाद्यैश्च पुराणपठनादिभिः}
{तर्पयेद्गणनाथं च विप्रान् दानेन श्रोत्रियान्} %॥२४॥

\twolineshloka
{बन्धुभिः स्वजनैः सार्धं भुञ्जीयात्तैलवर्जितम्}
{एवं यः कुरुते मर्त्यो गणनाथप्रसादतः} %॥२५॥

\twolineshloka
{सिध्यन्ति सर्वकार्याणि नात्र कार्या विचारणा}
{ततः प्रभाते विमले पुनः पूजां समाचरेत्} %॥२६॥

\twolineshloka
{मौञ्जीं कृष्णाजिनं दण्डम् उपवीतं कमण्डलुम्}
{परिधानं तथा दद्याद्यथाविभवमुत्तमम्} %॥२७॥

\twolineshloka
{उपायनं ततो दद्यादाचार्याय स्वशक्तितः}
{अन्येभ्यो दक्षिणां दद्याद् ब्राह्मणान् भोजयेत्ततः} %॥२८॥

\twolineshloka
{त्रैलोक्यविश्रुतं चैतद् व्रतानामुत्तमोत्तमम्}
{अन्यैश्च देवमुनिभिर्गन्धर्वैः किन्नरैस्तथा} %॥२९॥

\twolineshloka
{चीर्णमेतत् सर्वैस्तु पुराकल्पे षडानन}
{इति पुत्राय शर्वेण षण्मुखायोदितं पुरा} %॥३०॥

\twolineshloka
{एवं कुरुष्व धर्मज्ञ गणनाथप्रपूजनम्}
{विजयस्ते भवेन्नित्यं सत्यं सत्यं वदाम्यहम्} %॥३१॥

\twolineshloka
{एतद् व्रतं हरिश्चापि दमयन्ती पुराऽकरोत्}
{कृष्णो जाम्बवतीमागाद् रत्नं चापि स्यमन्तकम्} %॥३२॥

\twolineshloka
{दमयन्ती नलं चैव व्रतस्यास्य प्रभावतः}
{शक्रेण पूजितः पूर्वं वृत्रासुरवधे तथा} %॥३३॥

\twolineshloka
{रामदेवेन तद्वच्च सीताया मार्गणे तथा}
{भगीरथेन तद्वच्च गङ्गामानयता पुरा} %॥३४॥

\twolineshloka
{अमृतोत्पादनार्थाय तथा देवासुरैरपि}
{कुष्ठव्याधियुतेनापि साम्बेनाराधितः पुरा} %॥३५॥

\twolineshloka
{एवमुक्तस्तु सूतेन सानुजः पाण्डुनन्दनः}
{पूजयामास देवस्य पुत्रं त्रिपुरघातिनः} %॥३६॥

\twolineshloka
{शत्रून् निहत्य प्राप्तवान् राज्यमोजसा}
{पूजयित्वा महाभागं गणेशं सिद्धिदायकम्} %॥३७॥

\twolineshloka
{सिध्यन्ति सर्वकार्याणि मनसा चिन्तितान्यपि}
{तेन ख्यातिं गतो लोके नाम्ना सिद्धिविनायकः} %॥३८॥

\twolineshloka
{विद्यारम्भे पूजितश्चेत् विद्यालाभो भवेद्ध्रुवम्}
{जयं च जयकामश्च पुत्रार्थी लभते सुतान्} %॥३९॥

\twolineshloka
{पतिकामा च भर्तारं सौभाग्यं च सुवासिनी}
{विधवा पूजयित्वा तु वैधव्यं नाप्नुयात्क्वचित्} %॥४०॥

\twolineshloka
{ब्राह्मणः क्षत्रियो वैश्यः शूद्रो वाऽप्यथवा स्त्रियः}
{अर्भकश्चापि भक्त्या च व्रतं कुर्याद्यथाविधि} %॥४१॥

\twolineshloka
{सिध्यन्ति सर्वकार्याणि गणनाथप्रसादतः}
{पुत्रपौत्राभिवृद्धिं च गजाद्यैश्वर्यमाप्नुयात्} %॥४२॥

॥ इति श्रीव्रतरत्नाकरे विनायकव्रतकल्पः सम्पूर्णः ॥
